% Chương 8 -- Hiện thực hệ thống, được đặt sau Sections/08-CaiDatHeThong.tex
% trong main.tex. Tên tệp cũ được giữ lại để hạn chế thay đổi đường dẫn.

\section{Hiện thực hệ thống}
\label{sec:hien-thuc-he-thong}

Chương này trình bày kết quả hiện thực giao diện web của hệ thống thương mại điện tử và trợ lý
mua sắm đa phương thức. Nội dung tập trung vào các màn hình mà người dùng trực tiếp tương tác,
cách các màn hình liên kết thành luồng nghiệp vụ, và nhóm chức năng quản trị phục vụ vận hành hệ
thống. Các quyết định kiến trúc, mô hình dữ liệu và giao tiếp dịch vụ đã được trình bày ở Chương
Thiết kế hệ thống; phần này chỉ mô tả kết quả hiện thực tương ứng trên Frontend.

% Chuẩn hóa kích thước các ảnh chụp giao diện trong Images/Frontend.
\newcommand{\frontendfigure}[3]{%
    \begin{figure}[H]
        \centering
        \IfFileExists{#1}{%
            \includegraphics[width=0.96\textwidth,height=0.72\textheight,keepaspectratio]{#1}%
        }{%
            \fbox{%
                \begin{minipage}[c][0.20\textheight][c]{0.90\textwidth}
                    \centering
                    Ảnh giao diện sẽ được bổ sung tại:\\[0.5em]
                    \texttt{\detokenize{#1}}
                \end{minipage}%
            }%
        }
        \caption{#2}
        \label{#3}
    \end{figure}%
}

\subsection{Tổng quan hiện thực Frontend}

Frontend được hiện thực bằng React và tổ chức theo bốn lớp trách nhiệm. Lớp định tuyến ánh xạ URL
đến từng trang chức năng; lớp trang điều phối dữ liệu và trạng thái của một màn hình; các component
dùng lại đảm nhiệm thẻ sản phẩm, bộ lọc, trạng thái rỗng, hộp thoại xem nhanh và tiến trình đơn
hàng; lớp API đóng gói giao tiếp với Backend thương mại điện tử và AI Service. Khu vực quản trị
được bảo vệ bằng tuyến riêng, trong khi các chức năng duyệt sản phẩm, tư vấn AI và đọc bài viết có
thể được truy cập từ tuyến công khai.

\begin{table}[H]
    \centering
    \caption{Ánh xạ các nhóm giao diện Frontend với chức năng chính}
    \label{tab:frontend-interface-map}
    \begin{tabularx}{\textwidth}{|>{\RaggedRight\arraybackslash}p{0.22\textwidth}|>{\RaggedRight\arraybackslash}p{0.25\textwidth}|X|}
        \hline
        \textbf{Nhóm giao diện} & \textbf{Tuyến tiêu biểu} & \textbf{Chức năng hiện thực} \\
        \hline
        Khám phá sản phẩm & \texttt{/}, \texttt{/category} & Hiển thị danh mục, danh sách sản phẩm, tìm kiếm, lọc và điều hướng đến chi tiết sản phẩm. \\
        \hline
        Hỗ trợ quyết định & \texttt{/prodetail/:id}, \texttt{/compare}, \texttt{/wishlist} & Xem thông tin chi tiết, so sánh sản phẩm, lưu danh sách yêu thích và xem nhanh sản phẩm. \\
        \hline
        Trợ lý AI & \texttt{/tu-van}, \texttt{/ai} & Hội thoại nhiều lượt, gửi văn bản hoặc hình ảnh và hiển thị sản phẩm được đề xuất. \\
        \hline
        Mua hàng & \texttt{/cart}, \texttt{/checkout} & Quản lý giỏ hàng, nhập thông tin nhận hàng, chọn phương thức thanh toán và xác nhận đơn. \\
        \hline
        Tài khoản & \texttt{/login}, \texttt{/register}, \texttt{/info}, \texttt{/order-history} & Xác thực, cập nhật hồ sơ, quản lý địa chỉ và theo dõi lịch sử đơn hàng. \\
        \hline
        Quản trị & \texttt{/admin} & Theo dõi tổng quan và quản lý sản phẩm, khách hàng, đơn hàng, khuyến mãi, coupon và bài viết. \\
        \hline
    \end{tabularx}
\end{table}

\subsection{Giao diện dành cho khách hàng}

\subsubsection{Trang chủ và danh mục sản phẩm}

Trang chủ là điểm bắt đầu của hành trình mua sắm, kết hợp banner, nhóm danh mục, sản phẩm nổi bật
và các nội dung hỗ trợ khám phá. Từ đây, người dùng chuyển sang trang danh mục để tìm kiếm, lọc và
duyệt danh sách sản phẩm. Cách tổ chức này giữ thao tác khám phá ở một luồng liên tục, đồng thời
tận dụng các component thẻ sản phẩm dùng chung để bảo đảm cách trình bày nhất quán.

\frontendfigure{Images/Frontend/giao diện trang chủ của hệ thống.png}{Giao diện trang chủ của hệ thống}{fig:frontend-home}

\frontendfigure{Images/Frontend/giao diện danh mục, tìm kiếm và lọc sản phẩm.png}{Giao diện danh mục, tìm kiếm và lọc sản phẩm}{fig:frontend-category}

\subsubsection{Chi tiết và hỗ trợ lựa chọn sản phẩm}

Trang chi tiết tổng hợp hình ảnh, giá, biến thể, trạng thái tồn kho và đánh giá để người dùng có đủ
thông tin trước khi thêm sản phẩm vào giỏ. Bên cạnh đó, chức năng so sánh, danh sách yêu thích, sản
phẩm đã xem gần đây và hộp thoại xem nhanh giúp người dùng rút ngắn quá trình cân nhắc mà không làm
mất ngữ cảnh duyệt hiện tại.

\frontendfigure{Images/Frontend/giao diện chi tiết sản phẩm.png}{Giao diện chi tiết sản phẩm}{fig:frontend-product-detail}

Trang sản phẩm yêu thích tập hợp những mặt hàng người dùng đã lưu trong quá trình duyệt. Từ danh
sách này, người dùng có thể quay lại trang chi tiết, thêm sản phẩm phù hợp vào giỏ hoặc loại bỏ
những lựa chọn không còn quan tâm.

\frontendfigure{Images/Frontend/giao diện trang sản phẩm yêu thích.png}{Giao diện danh sách sản phẩm yêu thích}{fig:frontend-wishlist}

\subsubsection{Tư vấn sản phẩm bằng AI}

Hệ thống cung cấp hai hình thức tương tác hỗ trợ mua sắm. Chatbot nổi cho phép đặt câu hỏi nhanh
ngay trong lúc duyệt các trang thương mại điện tử. Trang trợ lý chuyên biệt cung cấp không gian lớn
hơn cho hội thoại nhiều lượt, gửi văn bản hoặc hình ảnh tham chiếu và xem các sản phẩm được đề xuất.
Kết quả được đặt trong ngữ cảnh câu trả lời để người dùng có thể mở trang chi tiết và tiếp tục quy
trình mua sắm.

\frontendfigure{Images/Frontend/giao diện chatbot.png}{Giao diện chatbot hỗ trợ người dùng trong quá trình mua sắm}{fig:frontend-chatbot}

\frontendfigure{Images/Frontend/giao diện trợ lý mua sắm đa phương tiện.png}{Giao diện trợ lý mua sắm đa phương thức}{fig:frontend-ai-consulting}

\subsubsection{Giỏ hàng và thanh toán}

Giỏ hàng cho phép cập nhật số lượng, loại bỏ mặt hàng và quan sát lại tổng tiền trước khi chuyển sang
checkout. Tại bước thanh toán, hệ thống thu thập địa chỉ nhận hàng, áp dụng mã giảm giá hợp lệ và cho
phép lựa chọn phương thức thanh toán. Thông tin tổng kết đơn hàng được duy trì xuyên suốt hai màn
hình để người dùng có thể kiểm tra quyết định trước khi xác nhận.

\frontendfigure{Images/Frontend/giao diện quản lý giỏ hàng.png}{Giao diện quản lý giỏ hàng}{fig:frontend-cart}

\frontendfigure{Images/Frontend/giao diện thanh toán.png}{Giao diện checkout và lựa chọn phương thức thanh toán}{fig:frontend-checkout}

Sau khi đơn được ghi nhận thành công, hệ thống hiển thị màn hình xác nhận để người dùng nhận biết
quy trình đặt hàng đã hoàn tất và có thể chuyển đến lịch sử đơn hàng để tiếp tục theo dõi.

\frontendfigure{Images/Frontend/giao diện đặt hàng thành công.png}{Giao diện xác nhận đặt hàng thành công}{fig:frontend-order-success}

\subsubsection{Đăng nhập và đăng ký tài khoản}

Giao diện đăng nhập cho phép khách hàng truy cập các chức năng cá nhân hóa và tiếp tục hành trình
mua sắm bằng tài khoản đã có. Với người dùng mới, giao diện đăng ký thu thập các thông tin cần thiết,
kiểm tra dữ liệu nhập và tạo tài khoản trước khi chuyển sang bước xác thực hoặc đăng nhập hệ thống.

\frontendfigure{Images/Frontend/giao diện đăng nhập.png}{Giao diện đăng nhập hệ thống}{fig:frontend-login}

\frontendfigure{Images/Frontend/giao diện đăng ký.png}{Giao diện đăng ký tài khoản}{fig:frontend-register}

\subsubsection{Tài khoản và theo dõi đơn hàng}

Sau khi đăng nhập, người dùng có thể quản lý hồ sơ cá nhân và theo dõi lịch sử đơn hàng. Trạng thái
xử lý được trình bày theo dòng thời gian; người dùng có thể xem lại các mặt hàng đã mua và gửi phản
hồi khi phù hợp. Các trang cần phân biệt rõ trạng thái đang tải, không có dữ liệu và lỗi tải dữ liệu
để tránh tạo cảm giác hệ thống không phản hồi.

\frontendfigure{Images/Frontend/giao diện tài khoản.png}{Giao diện quản lý thông tin tài khoản}{fig:frontend-account}

\frontendfigure{Images/Frontend/giao diện lịch sử đơn hàng.png}{Giao diện lịch sử và trạng thái đơn hàng}{fig:frontend-order-history}

\subsection{Giao diện dành cho quản trị viên}

Khu vực quản trị sử dụng bố cục riêng với thanh điều hướng bên, chỉ cho phép tài khoản có quyền phù
hợp truy cập. Trang tổng quan cung cấp các chỉ số vận hành cần theo dõi; các màn hình nghiệp vụ cho
phép quản lý sản phẩm, khách hàng, đơn hàng, khuyến mãi, coupon và bài viết. Những thao tác tạo, sửa
hoặc chuyển trạng thái được thực hiện qua biểu mẫu và hộp thoại xác nhận nhằm hạn chế thay đổi ngoài
ý muốn.

\frontendfigure{Images/Frontend/giao diện tổng quan dành cho quản trị viên.png}{Giao diện tổng quan dành cho quản trị viên}{fig:frontend-admin-dashboard}

\subsubsection{Quản lý sản phẩm, người dùng và đơn hàng}

Các màn hình quản lý dữ liệu cốt lõi sử dụng bảng để quản trị viên quan sát nhiều bản ghi, tìm kiếm
và thực hiện thao tác nghiệp vụ tại cùng một nơi. Quản lý sản phẩm hỗ trợ cập nhật thông tin hàng
hóa; quản lý người dùng phục vụ kiểm soát tài khoản; quản lý đơn hàng cho phép theo dõi và chuyển
trạng thái phù hợp với vòng đời đã thiết kế.

\frontendfigure{Images/Frontend/giao diện quản lý sản phẩm.png}{Giao diện quản lý sản phẩm}{fig:frontend-admin-products}

\frontendfigure{Images/Frontend/giao diện quản lý người dùng.png}{Giao diện quản lý người dùng}{fig:frontend-admin-users}

\frontendfigure{Images/Frontend/giao diện quản lý đơn hàng.png}{Giao diện quản lý đơn hàng}{fig:frontend-admin-orders}

\subsubsection{Quản lý khuyến mãi và mã giảm giá}

Hai màn hình khuyến mãi và mã giảm giá tách biệt chính sách giảm giá tự động khỏi mã do khách hàng
nhập. Quản trị viên có thể theo dõi thời gian hiệu lực, trạng thái và phạm vi áp dụng của từng cấu
hình, qua đó giảm nhầm lẫn khi nhiều chương trình được vận hành đồng thời.

\frontendfigure{Images/Frontend/giao diện quản lý khuyến mãi.png}{Giao diện quản lý chương trình khuyến mãi}{fig:frontend-admin-promotions}

\frontendfigure{Images/Frontend/giao diện quản lý mã giảm giá.png}{Giao diện quản lý mã giảm giá}{fig:frontend-admin-coupons}

\subsubsection{Quản lý bài viết}

Giao diện quản lý bài viết hỗ trợ quản trị nội dung giới thiệu, hướng dẫn và tin tức của cửa hàng.
Việc tách nội dung này khỏi danh mục sản phẩm giúp vòng đời biên tập không phụ thuộc vào các thao
tác cập nhật hàng hóa.

\frontendfigure{Images/Frontend/giao diện quản lý bài viết.png}{Giao diện quản lý bài viết}{fig:frontend-admin-blogs}

\subsection{Nguyên tắc phản hồi và khả năng sử dụng}

Các giao diện được hiện thực theo một tập nguyên tắc thống nhất: phản hồi ngay khi người dùng thực
hiện thao tác; khóa nút gửi trong thời gian xử lý để hạn chế yêu cầu trùng; dùng skeleton cho nội dung
đang tải; hiển thị trạng thái rỗng có hướng dẫn hành động; và thông báo lỗi bằng ngôn ngữ gắn với tác
vụ đang thực hiện. Những trang chứa danh sách dài sử dụng phân trang hoặc lọc, còn các thao tác ảnh
hưởng đến đơn hàng và dữ liệu quản trị phải có bước xác nhận rõ ràng.

Giao diện được tổ chức theo component và CSS module để giới hạn phạm vi kiểu trình bày, đồng thời
cho phép tái sử dụng các thành phần ở nhiều trang. Việc kiểm tra trên các kích thước màn hình khác
nhau cần bảo đảm thanh điều hướng, biểu mẫu, bảng dữ liệu và hộp thoại vẫn sử dụng được mà không che
khuất thao tác chính.

\subsection*{Kết chương}

Chương này đã trình bày kết quả hiện thực Frontend theo hai phạm vi khách hàng và quản trị viên.
Các giao diện bao phủ hành trình từ khám phá, nhận tư vấn, lựa chọn sản phẩm, đặt hàng đến theo dõi
đơn; đồng thời cung cấp công cụ vận hành dữ liệu cho quản trị viên. Chương tiếp theo trình bày kế
hoạch kiểm thử và kết quả đánh giá các chức năng AI cũng như toàn bộ hệ thống.
